% !TeX root = ../main.tex
\section{引言}
\label{sec:introduction}
%%下面的文字内容，仅支持小范围的文本描述修改优化，不允许大范围重写
随着大模型技术的发展，模型的参数规模与计算需求持续增长，采用更低精度的数值格式已成为常见的优化手段~\cite{deepseekv3}。通用矩阵乘法（GEMM）作为大模型的主要计算负载，其计算格式也逐渐由单精度浮点格式（FP32）向低精度格式拓展。其中，最主要的代表包括 TensorFloat-32（TF32）、半精度浮点格式（FP16）和 BF16。随着大模型量化技术的发展，FP8、INT8、FP6、FP4，以及引入共享缩放因子（scaling factor）的微缩放（Microscaling）格式~\cite{ocpmx2023}，也逐渐成为模型推理优化中的重要选择。以 DeepSeek-V4.1-Flash~\cite{deepseekv41source}为例，其文本主干采用按算子配置的混合精度策略：混合专家（MoE）计算及多数注意力投影采用 FP8；稀疏注意力索引器的 QK 路径采用 MXFP4；主注意力的 QK/PV 计算采用 BF16。因此可见，当前大语言模型推理呈现出了多种精度计算需求并存的特点，对承载大模型推理的硬件也带来了更高的要求。

在硬件设计方面，在相同的数据输入位宽下，低精度格式可同时传入更多元素，形成BF16:MXFP8:MXFP4之间$1:2:4$的峰值算力配比。但是，当前工业界的最新AI芯片也逐渐打破这种传统的模式，通过定向提高FP4算力，或同时增强FP8与FP4的峰值算力，来增强产品竞争力。但是模型的发展的节奏与芯片设计制造周期的之间有较大差异,根据当前的模型需求设计芯片可能无法满足未来的发展变化。因此芯片设计必须为未来的模型负载变化留有设计余量。从AI芯片的主要算力来源——张量计算单元（Tensorcore）来分析，需要在支持多种计算精度的基础上，还需要对不同精度之间的相对算力进行取舍，留有计算容量，兼顾当前与未来的下模型发展需求，并权衡由此带来的硬件成本和收益。

因此针对模型负责变化选择合适的芯片算力配比，这一极具工程价值和实际意义的问题，我们提出了DP-Puzzle ,将算力配比选择的硬件成本与模型负载收益联系起来。由于Tensorcore的设计中，通常采用点积单元结合脉动~\cite{jouppi2017tpu}或者广播的数据传输格式，构建二维计算阵列。因此针对算力配比问题，我们以点积单元为主要的算力配比改变的硬件设计实现，支持多种数据精度格式的计算，支持改变低精度的算力。根据真实的面积功耗硬件成本，作为算力配比分析的基础。结合硬件成本和模型负载收益，分析不同模型负载迁移下的不同算力配比的综合收益，为芯片设计配比的规格定义提供参考依据。

% BEGIN INLINE TABLE: related_features
% Promoted from the approved temporary manuscript.
\suppressfloats[t]
\begin{table}[t]
\centering
\caption{精度支持与组合式配比配置能力。}
\label{tab:related_features}
\begingroup\footnotesize\setlength{\tabcolsep}{1pt}
\def\Y{$\checkmark$}\def\N{$\times$}\def\E{$\triangle$}
\begin{tabular*}{\columnwidth}{@{\extracolsep{\fill}}l*{10}{c}@{}}\toprule
Design & TF32 & FP16 & BF16 & \shortstack{MX\\FP8} & \shortstack{MX\\FP6} & \shortstack{MX\\FP4} & INT8 & \shortstack{NV\\FP4} & HiF4 & \shortstack{Multi\\Ratio} \\\midrule
Ten-Four~\cite{rout2026tenfourv3} & \Y & \Y & \Y & \Y & \N & \N & \Y & \N & \N & \N \\
MXDOTP~\cite{islamoglu2025mxdotp} & \N & \N & \N & \Y & \N & \N & \N & \N & \N & \N \\
FPnew~\cite{mach2021fpnew} & \N & \Y & \Y & \E & \N & \N & \N & \N & \N & \N \\
MiniFloats~\cite{bertaccini2024minifloats} & \N & \Y & \Y & \E & \N & \N & \N & \N & \N & \N \\
TransDot~\cite{transdot2026} & \N & \Y & \N & \E & \N & \E & \N & \N & \N & \N \\
\midrule
DP-Puzzle & \Y & \Y & \Y & \Y & \Y & \Y & \Y & \Y & \Y & \Y \\
\bottomrule\end{tabular*}
\par\smallskip
\parbox{\columnwidth}{\raggedright
$\checkmark$: supported; $\times$: absent in the reported configuration; $\triangle$: element arithmetic supported, integrated MX scaling not established. FPnew and MiniFloats also support FP32/FP64; MiniFloats uses FP8 E5M2/E4M3 and BF16-encoded FP16alt.}
\endgroup
\end{table}

% END INLINE TABLE: related_features
本文的主要贡献如下：
\begin{enumerate}
\item \textbf{可组合的多精度点积微架构。} 我们提出 DP-Puzzle，采用 Lane、Group 和点积单元三级架构，支持 TF32、FP16/BF16、INT8、MXFP8/6/4、NVFP4 和 HiF4。设计在不同精度间复用乘法、对齐与规约资源，并通过增配专用 Group 独立扩展低精度吞吐。
\item \textbf{九种算力配比的实现与评估。} 我们实现九种 BF16:MXFP8:MXFP4 峰值算力配比，覆盖基线、三种产业配比及五种扩展配置，并基于 12\,nm 逻辑综合与负载驱动功耗分析量化其实现成本。相较于 \mbox{$1:2:4$}，\mbox{$1:2:8$} 将 MXFP4 吞吐翻倍，面积增加 33.60\%，MXFP4 面效和能效分别提高 49.70\% 和 27.27\%。
\item \textbf{联合预算下的配比选择。} 我们在相同总核心面积和点积阵列功耗预算下，估算各配比的可部署核心数与分精度峰值算力。在 DeepSeek-V4.1-Flash prefill 的五种精度迁移场景和 11 个上下文长度的综合评估中，\mbox{$1:2:8$} 在九种配比中表现最佳，相对 \mbox{$1:2:4$} 取得 $1.13\times$ 的理想 GEMM 计算加速。
\end{enumerate}

本文其余部分安排如下：第~\ref{sec:background}~节介绍背景与动机；第~\ref{sec:architecture}~节阐述组合式点积微架构与配比构造；第~\ref{sec:evaluation}~节给出评估方法及面积、功耗结果；第~\ref{sec:workload}~节分析面积、功耗约束与精度负载迁移下的配比选择；第~\ref{sec:conclusion}~节总结全文。
